% ============================================================================
%  section_sm_identity_layer.tex
%  VSEPR-SIM --- SM-Channel Identity Layer
%  Theory, mathematics, and code reference for the WO-SM-IDENTITY overhaul.
%
%  Compile:  pdflatex section_sm_identity_layer.tex  (run twice for TOC)
%  v5.2.0  |  WO-SM-IDENTITY  |  v5.0.0-main
% ============================================================================
\documentclass[11pt,a4paper]{article}

\usepackage[margin=1in]{geometry}
\usepackage{amsmath,amssymb,amsthm}
\usepackage{booktabs}
\usepackage{longtable}
\usepackage{xcolor}
\usepackage{listings}
\usepackage{hyperref}
\usepackage{titlesec}
\usepackage[protrusion=true,expansion=false]{microtype}
\usepackage{parskip}
\usepackage{mathtools}
\usepackage{array}

% ---------------------------------------------------------------------------
%  Colour palette
% ---------------------------------------------------------------------------
\definecolor{codebg}{RGB}{245,246,248}
\definecolor{kw}{RGB}{0,80,160}
\definecolor{cm}{RGB}{60,128,60}
\definecolor{st}{RGB}{160,40,40}
\definecolor{nu}{RGB}{140,60,180}
\definecolor{pp}{RGB}{100,50,0}
\definecolor{rulecol}{RGB}{180,190,200}
\definecolor{boxbg}{RGB}{240,248,255}

% ---------------------------------------------------------------------------
%  C++23 listing style
% ---------------------------------------------------------------------------
\lstdefinestyle{cpp23}{
  language=C++,
  backgroundcolor=\color{codebg},
  basicstyle=\ttfamily\small,
  keywordstyle=\color{kw}\bfseries,
  commentstyle=\color{cm}\itshape,
  stringstyle=\color{st},
  numberstyle=\color{nu}\tiny,
  emphstyle=\color{pp}\bfseries,
  emph={ParticleIdentity,InteractionCandidate,InteractionChannelRow,
		ConservationGates,ChannelPolicy,KappaScores,Matrix2x2,Matrix3x3,
		SpinProxy,ColorCharge,QuarkFlavor,evaluate_pair,compute_kappa,
		kappa_total,score_channels,evaluate_conservation,
		qcd_accumulate_forces_sm,double,int64_t,uint64_t,int32_t,int8_t,bool},
  breaklines=true,
  frame=single,
  rulecolor=\color{rulecol},
  framesep=4pt,
  numbers=left,
  numbersep=8pt,
  tabsize=4,
  showstringspaces=false,
  xleftmargin=14pt,
  captionpos=b
}
\lstset{style=cpp23}

% ---------------------------------------------------------------------------
%  Section formatting
% ---------------------------------------------------------------------------
\titleformat{\section}{\large\bfseries}{}{0em}{}[\titlerule]
\titleformat{\subsection}{\normalsize\bfseries}{}{0em}{}

% ---------------------------------------------------------------------------
%  Theorem-like environments
% ---------------------------------------------------------------------------
\newtheoremstyle{defbox}
  {6pt}{6pt}{\normalfont}{0pt}{\bfseries}{.}{0.5em}{}
\theoremstyle{defbox}
\newtheorem{definition}{Definition}
\newtheorem{principle}{Principle}

% ---------------------------------------------------------------------------
\begin{document}

\begin{center}
  {\LARGE\bfseries SM-Channel Identity Layer}\\[4pt]
  {\large VSEPR-SIM --- Theory, Mathematics, and Code Reference}\\[2pt]
  {\small\texttt{include/IDENTITY/} \quad \texttt{include/PARTICLES/}
		 \quad \texttt{include/COLLISION/}}\\[6pt]
  {\small v5.2.0 \quad WO-SM-IDENTITY \quad v5.0.0-main}
\end{center}

\vspace{4pt}
\hrule
\vspace{10pt}

\tableofcontents
\newpage

% ============================================================================
\section{Motivation and Design Philosophy}
\label{sec:motivation}
% ============================================================================

Classical particle simulations traditionally compute interactions as a single
universal force rule evaluated over every pair.  A scalar force magnitude is
produced, integrated, and repeated.  This correctly models billiard-ball
mechanics and is inadequate for anything more interesting.

The Standard Model does not work this way.  It treats particle interactions as
\emph{allowed transitions between field states}, constrained by conservation
laws, quantum number compatibility, and the discrete set of force carriers
available for a given pair.  A photon is exchanged, or a $W$ boson is, or a
gluon is.  Which carrier mediates the interaction determines everything:
whether identity is preserved, mutated, branched, confined, or annihilated.

The SM-channel identity layer introduces this logic into VSEPR-SIM.  It does
not replace the Cornell QCD force kernel.  It wraps around it as a gate layer,
selecting which physics fires per pair and recording what identity information
is gained or lost in the process.

\begin{principle}[Channel Selection]
Particle-particle interactions are modeled as gated transitions between
recoverable identity-information states.  Standard Model-inspired channels
determine whether identity is preserved, transformed, branched, confined,
or annihilated.
	\end{principle}

\noindent\textbf{IKK theoretical grounding.}
The Inverse Kaluza-Klein (IKK) framework (\texttt{IKK\_Doc\_I}, \texttt{IKK\_Doc\_II})
formalises the same intuition at a more fundamental level: the observed
wavefunction $\Psi(\mathbf{r},t)$ is a \emph{scale-projected shadow} of a fuller
identity structure $\mathcal{I}$ defined over a scale ladder
$\{S_0, S_1, \ldots, S_D\}$.  Failures of the standard equation at high
energy, many-body, and cosmological scales are \emph{projection losses} ---
physics present in $\mathcal{I}$ but not captured by the $S_3/S_4$ projection.
The SM-channel layer is the VSIM operationalisation of this idea:
each particle carries a \texttt{scale\_mask} bitmask naming which scale levels
are active, and hidden-intensity fields ($H_\text{charge}$, $H_\text{color}$,
$H_\text{spin}$) track structure invisible to the net-zero observable channel.
See \texttt{section\_ikk\_kernel\_bridge} for the full theoretical-to-code mapping.

% ============================================================================
\section{Notation and Definitions}
\label{sec:notation}
% ============================================================================

\subsection{Particles, Channels, and Output}

Let $C_i$ denote the \emph{complete state} of particle $i$: its position,
velocity, and full identity bundle.  Let $I_i$ denote the
\emph{recoverable identity-information} scalar for particle $i$.  An
interaction between particles $i$ and $j$ mediated by channel $A_{ij}$ is:

\begin{equation}
  C_i + C_j \xrightarrow{\,A_{ij}\,} I_{\mathrm{out}}
  \qquad \text{(scattering / identity-preserving)}
\end{equation}

\begin{equation}
  I_i + I_j \xrightarrow{\,A_{ij}\,} C_{\mathrm{products}}
  \qquad \text{(annihilation)}
\end{equation}

\subsection{Interaction Channel Row}

For a pair $(i,j)$, the full channel row $M_{ij}$ is:

\begin{equation}
  M_{ij} = \bigl[\,
	m_\gamma,\;
	m_g,\;
	m_W,\;
	m_Z,\;
	m_H,\;
	m_{\mathrm{ann}},\;
	m_{\mathrm{pair}},\;
	m_{\mathrm{decay}},\;
	m_{\mathrm{conf}}
  \,\bigr]_{ij}
\end{equation}

The aggregate interaction amplitude is:

\begin{equation}
  A_{ij} = \sum_k M_{ij}^{(k)}\,K_k\!\left(C_i, C_j, I_i, I_j, Z_i, Z_j\right)
\end{equation}

Table~\ref{tab:channels} defines each channel.

\begin{table}[ht]
\centering
\caption{SM-inspired interaction channels in $M_{ij}$.}
\label{tab:channels}
\renewcommand{\arraystretch}{1.25}
\begin{tabular}{llll}
\toprule
\textbf{Symbol} & \textbf{SM analogy} & \textbf{Status} & \textbf{What is logged} \\
\midrule
$m_\gamma$              & EM / photon-mediated       & \textbf{active}  & charge, scattering, radiation \\
$m_g$                   & Strong / gluon-mediated    & \textbf{active}  & color flow, Cornell gate \\
$m_W$                   & Weak charged current       & gated/zero       & flavor/charge-changing transition \\
$m_Z$                   & Weak neutral current       & gated/zero       & weak scattering, no charge exchange \\
$m_H$                   & Higgs / scalar coupling    & gated/zero       & mass-coupling interaction \\
$m_{\mathrm{ann}}$      & Annihilation               & \textbf{active}  & particle-antiparticle collapse score \\
$m_{\mathrm{pair}}$     & Pair production            & gated/zero       & energy-to-particle creation \\
$m_{\mathrm{decay}}$    & Unstable decay             & gated/zero       & lifetime, branching ratio \\
$m_{\mathrm{conf}}$     & Confinement/hadronization  & gated/zero       & no free color states \\
\bottomrule
\end{tabular}
\end{table}

\subsection{The $Z$-Matrix}

Each particle carries two identity matrices:

\begin{definition}[$Z$-matrix]
$Z_i^{(2)} \in \mathbb{R}^{2\times 2}$ is the \emph{runtime state matrix} of
particle $i$.  It encodes the particle's current field configuration and is
updated by the integrator after each interaction step.

$Z_i^{(3)} \in \mathbb{R}^{3\times 3}$ is the \emph{fundamental identity
matrix} in the three-generation space.  It changes only during weak-mediated
identity mutation events.
\end{definition}

Both matrices are plain-aggregate types storing \texttt{double[n][n]} in
row-major order.  The key operation is \emph{matrix opposition}:

\begin{equation}
  \kappa_Z(Z_i, Z_j)
  = 1 - \frac{\langle Z_i,\, Z_j \rangle_F}{\|Z_i\|_F\,\|Z_j\|_F + \varepsilon}
\end{equation}

where $\langle A, B\rangle_F = \sum_{r,c} A_{rc}B_{rc}$ is the Frobenius inner
product.  The range is $[0, 1]$: zero means the matrices are identical
(same state), one means they are fully opposed (antiparticle pairing).

An alternative form via the trace product is:

\begin{equation}
  \kappa_Z = 1 - \frac{\operatorname{tr}(Z_i Z_j)}
					  {\|Z_i\|_F\,\|Z_j\|_F + \varepsilon}
\end{equation}

\noindent\textbf{IKK connection: hidden-active condition.}
IKK Doc I~(\S{}3.3) and Doc II~(\S{}D) establish that a net-zero observable
channel ($I_j = \alpha^\top \mathbf{c}_j = 0$)
absence of internal structure.  The hidden intensity
$H_j = \mathbf{c}_j^\top \mathbf{W}\, \mathbf{c}_j > 0$ can coexist with
$I_j = 0$.  Canonical example: a neutral baryon (up+down+down quarks) has
$I_\text{charge} = +\tfrac{2}{3} - \tfrac{1}{3} - \tfrac{1}{3} = 0$
but $H_\text{charge} = \tfrac{2}{3} \neq 0$.
The \texttt{ParticleIdentity} fields \texttt{H\_charge}, \texttt{H\_color},
\texttt{H\_spin} and the \texttt{hidden\_active\_charge()} member implement
this detection.  The \texttt{Z3} $3\times 3$ matrix is the direct analog of
the quark component matrix $\mathbf{C}_q$ from IKK Doc II~(\S{}C).

% ============================================================================
\section{Conservation Gates}
\label{sec:gates}
% ============================================================================

Every interaction is conditional on a conservation ledger.  The gate product
is:

\begin{equation}
  \mathcal{G}_{\mathrm{cons}}
  = G_E \cdot G_p \cdot G_q \cdot G_\ell \cdot G_b \cdot G_c \cdot G_s
\end{equation}

Table~\ref{tab:gates} defines each gate.

\begin{table}[ht]
\centering
\caption{Conservation gates and their checks.}
\label{tab:gates}
\renewcommand{\arraystretch}{1.25}
\begin{tabular}{lll}
\toprule
\textbf{Gate} & \textbf{Quantity} & \textbf{Condition checked} \\
\midrule
$G_E$      & Energy conservation   & $E_{\mathrm{rel}} \geq 0$, kinematically allowed \\
$G_p$      & Momentum balance      & $\Delta p / p_{\mathrm{total}} < \epsilon_p$ \\
$G_q$      & Electric charge       & $\sum q$ is finite and consistent \\
$G_\ell$   & Lepton family         & families match or at least one is zero \\
$G_b$      & Baryon number         & additive, tracked per interaction \\
$G_c$      & Color neutrality      & colored pairs satisfy $c_j = \overline{c_i}$ \\
$G_s$      & Spin compatibility    & spin types are pairwise compatible \\
\bottomrule
\end{tabular}
\end{table}

The annihilation condition requires all gates plus spatial/kinematic gates:

\begin{equation}
  \boxed{
	A_{ij} =
	  G_{\mathrm{cons}} \cdot G_{\mathrm{kin}} \cdot G_{\mathrm{overlap}}
	  \cdot G_\kappa \cdot G_L
  }
\end{equation}

where:

\begin{align}
  G_{\mathrm{overlap}} &: r_{ij} < r_c \\
  G_{\mathrm{kin}}     &: E_{\mathrm{rel}} > E_c \\
  G_\kappa             &: \kappa_{ij}^{\mathrm{total}} > \kappa_c \\
  G_L                  &: L_{ij} > L_{\mathrm{ann}}
\end{align}

% ============================================================================
\section{Kappa: Compatibility and Opposition Scores}
\label{sec:kappa}
% ============================================================================

\subsection{Sub-scores}

Five orthogonal compatibility scores are computed per pair:

\begin{align}
  \kappa_Z &= 1 - \frac{\langle Z_i, Z_j\rangle_F}{\|Z_i\|_F\|Z_j\|_F+\varepsilon}
  \tag{matrix opposition} \\[4pt]
  \kappa_q &= 1 - \frac{|q_i + q_j|}{|q_i|+|q_j|+\varepsilon}
  \tag{charge complementarity} \\[4pt]
  \kappa_s &=
	\begin{cases} 1 & s_i = s_j \\ 0.5 & \text{otherwise} \end{cases}
  \tag{spin compatibility} \\[4pt]
  \kappa_c &=
	\begin{cases}
	  1 & c_j = \overline{c_i},\; \text{both colored} \\
	  0 & \text{otherwise}
	\end{cases}
  \tag{color/anticolor match} \\[4pt]
  \kappa_E &= 1
  \tag{energy availability (placeholder)}
\end{align}

\subsection{Weighted Total}

\begin{equation}
  \boxed{
	\kappa_{ij}^{\mathrm{total}}
	= w_Z\,\kappa_Z
	+ w_q\,\kappa_q
	+ w_s\,\kappa_s
	+ w_c\,\kappa_c
	+ w_E\,\kappa_E
  }
\end{equation}

Default weights in \texttt{ChannelPolicy}:

\begin{center}
\begin{tabular}{lrrrrr}
\toprule
Weight & $w_Z$ & $w_q$ & $w_s$ & $w_c$ & $w_E$ \\
\midrule
Value  & 0.40  & 0.30  & 0.10  & 0.10  & 0.10  \\
\bottomrule
\end{tabular}
\end{center}

% ============================================================================
\section{Identity Information and State Loss}
\label{sec:identity}
% ============================================================================

\subsection{Recoverable Identity Information}

Each particle carries a scalar $I_i \in [0, 1]$ representing how much of its
original identity information remains recoverable.  A freshly created particle
has $I_i = 1$.  After interactions:

\begin{definition}[State Loss Principle]
\begin{equation}
  \Delta I_{ij} = I_i^+ - I_i^-
  \qquad \text{(change over one interaction)}
\end{equation}
\begin{equation}
  L_{ij} = -\Delta I_{ij}
  \qquad \text{(state loss, } \geq 0 \text{ for ordinary loss)}
\end{equation}
\end{definition}

The three regimes:

\begin{center}
\renewcommand{\arraystretch}{1.3}
\begin{tabular}{lll}
\toprule
\textbf{Condition} & \textbf{Name} & \textbf{Meaning} \\
\midrule
$\Delta I_{ij} < 0$ & Ordinary loss       & Identity information is dissipated \\
$\Delta I_{ij} = 0$ & Identity-preserving & Elastic or neutral-current scattering \\
$\Delta I_{ij} > 0$ & Positive event      & Short-lived structured state forming \\
\bottomrule
\end{tabular}
\end{center}

\subsection{The Annihilation Threshold}

\begin{equation}
  G_L =
  \begin{cases}
	1 & L_{ij} > L_{\mathrm{ann}} \\
	0 & L_{ij} \leq L_{\mathrm{ann}}
  \end{cases}
  \qquad\text{where}\quad L_{ij} = -\Delta I_{ij}
\end{equation}

\subsection{Positive Information Events}

When $\Delta I_{ij} > 0$, a \emph{positive information event} is flagged.
This means:

\begin{equation}
  \boxed{
	\text{Recoverable identity temporarily becomes more structured before dissipating.}
  }
\end{equation}

Physical interpretations:

\begin{center}
\renewcommand{\arraystretch}{1.3}
\begin{tabular}{ll}
\toprule
\textbf{Event} & \textbf{Meaning} \\
\midrule
Temporary pairing       & Short-lived structured two-particle state \\
Resonance               & Intermediate state before decay \\
Matrix alignment        & $(Z_i, Z_j)$ become unusually compatible \\
Pre-annihilation order  & Final collapse is likely \\
Confinement ordering    & Colored components forming neutral products \\
Shower coherence        & Branching with organized angular/color structure \\
\bottomrule
\end{tabular}
\end{center}

% ============================================================================
\section{$\Delta I$ Heuristic Kernel}
\label{sec:delta_i}
% ============================================================================

The current $\Delta I$ computation is a heuristic pending a full
branching-ratio kernel.  It decomposes into four contributions:

\begin{align}
  \Delta I^{\mathrm{EM}}    &= -0.05\; m_\gamma
	\tag{small EM loss} \\[2pt]
  \Delta I^{\mathrm{ann}}   &= -0.90\; m_{\mathrm{ann}}
	\tag{near-total annihilation loss} \\[2pt]
  \Delta I^{\mathrm{color}} &= -0.20\; m_g\,\kappa_c
	\tag{color exchange loss} \\[2pt]
  \Delta I^{\mathrm{res}}   &= +0.10\,(1 - \kappa_Z)\;
							   \cdot\mathbf{1}[\,(1-\kappa_Z) > 0.80\,]
	\tag{resonance / alignment gain}
\end{align}

Total, clamped to $[-I_i,\; +0.15\,I_i]$:

\begin{equation}
  \Delta I_{ij}
  = \mathrm{clamp}\!\left(
	  \Delta I^{\mathrm{EM}}
	+ \Delta I^{\mathrm{ann}}
	+ \Delta I^{\mathrm{color}}
	+ \Delta I^{\mathrm{res}},\;
	-I_i,\;
	0.15\,I_i
	\right)
\end{equation}

% ============================================================================
\section{Full Annihilation Gate Stack}
\label{sec:annihilation}
% ============================================================================

\subsection{Electron-Positron Prototype}

The canonical annihilation prototype in VSEPR-SIM:

\begin{equation}
  e^- + e^+ \xrightarrow{A} \gamma_1 + \gamma_2
  \quad\Longleftrightarrow\quad
  I_{e^-} + I_{e^+} \xrightarrow{A} C_{\gamma_1} + C_{\gamma_2}
\end{equation}

Required conditions:

\begin{equation}
  q_i + q_j = 0
  \qquad
  p_i + p_j = p_{\mathrm{out}}
  \qquad
  E_i + E_j = E_{\mathrm{out}}
  \qquad
  \kappa_{ij}^{\mathrm{total}} > \kappa_c
\end{equation}

\subsection{Gate Stack Summary}

\begin{equation}
  \text{\texttt{annihilation\_candidate}} =
  \underbrace{r < r_c}_{G_{\mathrm{overlap}}}
  \;\land\;
  \underbrace{E_{\mathrm{rel}} > E_c}_{G_{\mathrm{kin}}}
  \;\land\;
  \underbrace{\kappa^{\mathrm{total}} > \kappa_c}_{G_\kappa}
  \;\land\;
  \underbrace{L_{ij} > L_{\mathrm{ann}}}_{G_L}
  \;\land\;
  \underbrace{\mathcal{G}_{\mathrm{cons}}}_{G_{\mathrm{cons}}}
\end{equation}

\begin{equation}
  \text{\texttt{annihilation\_confirmed}} =
  \text{\texttt{candidate}} \;\land\; |q_i + q_j| < 10^{-6}
\end{equation}

Default threshold values in \texttt{ChannelPolicy}:

\begin{center}
\begin{tabular}{llll}
\toprule
\textbf{Parameter} & \textbf{Field} & \textbf{Default} & \textbf{Meaning} \\
\midrule
$r_c$              & \texttt{r\_c}     & 0.20 & max separation (\AA\ analog) \\
$E_c$              & \texttt{E\_c}     & 0.50 & min relative energy \\
$\kappa_c$         & \texttt{kappa\_c} & 0.70 & min $\kappa^{\mathrm{total}}$ \\
$L_{\mathrm{ann}}$ & \texttt{L\_ann}   & 0.50 & min state loss \\
\bottomrule
\end{tabular}
\end{center}

% ============================================================================
\section{Strong Interaction and the Color Channel Gate}
\label{sec:qcd_gate}
% ============================================================================

The Cornell QCD kernel fires only when the color channel $m_g$ exceeds
the activation threshold $\delta_g$:

\begin{equation}
  \texttt{fire\_cornell} =
  \begin{cases}
	1 & m_g > \delta_g \;\;\text{(both particles color-charged)} \\
	0 & \text{otherwise}
  \end{cases}
  \qquad \delta_g = 0.05
\end{equation}

\subsection{Scale-Dependent Treatment}

\begin{center}
\renewcommand{\arraystretch}{1.25}
\begin{tabular}{ll}
\toprule
\textbf{Scale} & \textbf{Treatment} \\
\midrule
Very near       & Direct color/matrix interaction via $\kappa_c$ \\
Intermediate    & Cornell string/flux-tube proxy \\
Far             & Color-neutral aggregate only \\
Post-collision  & Shower + hadronization event model \\
\bottomrule
\end{tabular}
\end{center}

\subsection{Octree Color Neutrality Requirement}

\begin{equation}
  \boxed{
	\text{QCD octree aggregation must track color neutrality,
		  not just charge and mass.}
  }
\end{equation}

The planned \texttt{OctreeIdentityNode} adds to the existing Barnes--Hut
tree:
\begin{itemize}
  \item \texttt{color\_neutrality\_score} --- fraction of color-neutral aggregate
  \item \texttt{confinement\_risk} --- probability of forced confinement in
		one $T_B$ step
  \item \texttt{shower\_probability} --- probability of parton-shower initiation
\end{itemize}

% ============================================================================
\section{Weak Interaction and Identity Mutation}
\label{sec:weak}
% ============================================================================

Weak charged-current interactions can change particle type.  This is the
clean justification for an event where:

\begin{equation}
  I_i^- \neq I_i^+
  \qquad \text{(identity mutation, not just identity loss)}
\end{equation}

The beta-decay prototype in VSEPR-SIM notation:

\begin{equation}
  C_d \xrightarrow{A_W} C_u + C_{W^-}
  \qquad
  C_{W^-} \rightarrow C_{e^-} + C_{\bar\nu_e}
\end{equation}

The \texttt{weak\_identity\_transition} flag in \texttt{InteractionCandidate}
is set when $m_W > 0$.  The $W$ kernel is gated at zero in v5.2.0; the flag
exists for forward-compatibility with the live event logger.

% ============================================================================
\section{Code Reference: Data Structures}
\label{sec:code}
% ============================================================================

\subsection{Matrix Types (\texttt{IDENTITY/identity\_matrix.hpp})}

\begin{lstlisting}[caption={Matrix2x2 and Matrix3x3 plain aggregates.}]
struct Matrix2x2 {
	double m[2][2]{};
	double& operator()(int r, int c);
	static Matrix2x2 zero();
};

struct Matrix3x3 {
	double m[3][3]{};
	double& operator()(int r, int c);
	static Matrix3x3 zero();
};

// Free functions (same family for Matrix3x3):
Matrix2x2  identity_2x2();
double     trace(const Matrix2x2&);
double     frobenius_norm(const Matrix2x2&);
double     frobenius_dot(const Matrix2x2&, const Matrix2x2&);
double     opposition(const Matrix2x2&, const Matrix2x2&,
					  double eps = 1e-15);
double     matrix_compatibility(const Matrix2x2&, const Matrix2x2&,
								double eps = 1e-15);
Matrix2x2  mat_mul(const Matrix2x2&, const Matrix2x2&);
\end{lstlisting}

\subsection{ParticleIdentity (\texttt{IDENTITY/particle\_identity.hpp})}

\begin{lstlisting}[caption={ParticleIdentity struct.}]
struct ParticleIdentity {
	int64_t  id          {0};
	int32_t  type_code   {0};

	double    charge      {0.0};
	double    mass        {0.0};
	SpinProxy spin_proxy  {SpinProxy::Half};

	particles::ColorCharge  color      {ColorCharge::Neutral};
	particles::ColorCharge  anti_color {ColorCharge::Neutral};

	int8_t   lepton_family  {0};  // 0=none, 1=e, 2=mu, 3=tau
	int8_t   baryon_number  {0};
	int8_t   isospin_proxy  {0};

	Matrix2x2  Z2;   // runtime 2x2 state matrix
	Matrix3x3  Z3;   // 3x3 generation/flavor identity matrix

	double   I_recoverable  {1.0};
	uint64_t identity_hash  {0};

	void refresh_hash() noexcept;  // FNV-1a 64-bit over quantum-number fields
};

// Convenience constructors:
ParticleIdentity make_electron_like (int64_t id, double mass = 0.000511);
ParticleIdentity make_positron_like (int64_t id, double mass = 0.000511);
ParticleIdentity make_quark_like    (int64_t id, QuarkFlavor, ColorCharge);
\end{lstlisting}

\subsection{InteractionChannelRow (\texttt{IDENTITY/interaction\_channel.hpp})}

\begin{lstlisting}[caption={Nine-channel row and associated types.}]
struct InteractionChannelRow {
	double m_gamma {0.0};   // EM          (active)
	double m_g     {0.0};   // strong      (Cornell gate)
	double m_W     {0.0};   // weak charged  (gated/zero)
	double m_Z     {0.0};   // weak neutral  (gated/zero)
	double m_H     {0.0};   // Higgs         (gated/zero)
	double m_ann   {0.0};   // annihilation  (active)
	double m_pair  {0.0};   // pair prod.    (gated/zero)
	double m_decay {0.0};   // decay         (gated/zero)
	double m_conf  {0.0};   // confinement   (gated/zero)

	bool color_channel_active(double threshold = 0.05) const;
};

struct KappaScores {
	double kappa_Z, kappa_q, kappa_s, kappa_c, kappa_E;
};

struct ConservationGates {
	bool G_E, G_p, G_q, G_l, G_b, G_c, G_s;
	bool conservation_pass() const;  // AND of all gates
};

struct ChannelPolicy {
	double w_Z{0.40}, w_q{0.30}, w_s{0.10}, w_c{0.10}, w_E{0.10};
	double r_c{0.20}, E_c{0.50}, kappa_c{0.70}, L_ann{0.50};
	double m_g_threshold{0.05};
};
\end{lstlisting}

\subsection{InteractionCandidate (\texttt{IDENTITY/interaction\_candidate.hpp})}

\begin{lstlisting}[caption={Full pair evaluation record and evaluator.}]
struct InteractionCandidate {
	int64_t  i, j;               // particle indices
	double   r;                  // separation
	double   E_rel;              // relative kinetic energy
	double   p_balance_error;    // |delta_p| / |p_total|

	KappaScores           kappa;
	double                kappa_total_val;
	InteractionChannelRow channels;
	ConservationGates     gates;
	bool                  conservation_pass;

	double   delta_I;            // I_out - I_in
	double   state_loss;         // L_ij = -delta_I  (>= 0)

	double   annihilation_score;
	double   shower_score;
	double   weak_transition_score;
	double   resonance_score;

	bool     annihilation_candidate;
	bool     annihilation_confirmed;
	bool     positive_information_event;
	bool     weak_identity_transition;
	bool     color_exchange_active;
	bool     temporary_resonance;
};

// Main evaluator -- stateless, allocation-free.
InteractionCandidate evaluate_pair(
	const ParticleIdentity& a,
	const ParticleIdentity& b,
	const ChannelPolicy&    policy,
	double r, double E_rel,
	double p_balance_error = 0.0);
\end{lstlisting}

\subsection{Cornell Color-Channel Gate (\texttt{COLLISION/qcd\_force.hpp})}

\begin{lstlisting}[caption={\texttt{qcd\_accumulate\_forces\_sm()} -- SM-gated force pass.}]
// SM-gated overload.  Cornell fires only when m_g > policy.m_g_threshold.
// Original qcd_accumulate_forces() is untouched.
void qcd_accumulate_forces_sm(
	particles::QuarkGluonPlasma&                plasma,
	std::vector<identity::ParticleIdentity>&    identity_layer,
	const identity::ChannelPolicy&              policy,
	std::vector<identity::InteractionCandidate>* candidates_out);
\end{lstlisting}

% ============================================================================
\section{Planned Event Types for Live Logger}
\label{sec:events}
% ============================================================================

The live event logger (next pipeline stage) maps directly onto the
\texttt{InteractionCandidate} flag fields.

\begin{center}
\renewcommand{\arraystretch}{1.25}
\begin{tabular}{ll}
\toprule
\textbf{Event type string} & \textbf{Triggering condition} \\
\midrule
\texttt{elastic\_scattering}          & $\Delta I \approx 0$, no channel triggered \\
\texttt{inelastic\_scattering}        & $\Delta I < 0$, EM active \\
\texttt{radiative\_emission}          & $m_\gamma > 0$, significant loss \\
\texttt{pair\_production}             & $m_{\mathrm{pair}} > 0$ (future) \\
\texttt{annihilation\_candidate}      & \texttt{annihilation\_candidate == true} \\
\texttt{annihilation\_confirmed}      & \texttt{annihilation\_confirmed == true} \\
\texttt{weak\_identity\_transition}   & $m_W > 0$ (future) \\
\texttt{neutral\_current\_scattering} & $m_Z > 0$ (future) \\
\texttt{color\_exchange}              & \texttt{color\_exchange\_active == true} \\
\texttt{parton\_branch}               & shower initiated (future) \\
\texttt{hadronization\_candidate}     & $r > r_{\mathrm{conf}}$ (Cornell gate) \\
\texttt{temporary\_resonance}         & \texttt{temporary\_resonance == true} \\
\texttt{positive\_information\_event} & $\Delta I > 0$ \\
\texttt{identity\_loss\_spike}        & $L_{ij} \gg L_{\mathrm{ann}}$ \\
\bottomrule
\end{tabular}
\end{center}

\subsection{Example \texttt{.ee} Event Records}

\begin{verbatim}
frame 2091
event annihilation_candidate
particles 4001 4002
channel e_minus_e_plus_to_gamma_gamma
kappa_total 0.941
delta_I -0.882
state_loss 0.882
conservation_pass true
annihilation_score 0.918
end

frame 884
event temporary_resonance
particles 102 103
kappa_Z 0.031
delta_I +0.087
state_loss 0.000
positive_information_event true
end
\end{verbatim}

% ============================================================================
\section{Summary}
\label{sec:summary}
% ============================================================================

\begin{center}
\setlength{\fboxsep}{10pt}
\colorbox{boxbg}{%
\begin{minipage}{0.88\textwidth}
\centering
\textit{Particle-particle interactions are modeled as gated transitions between
recoverable identity-information states, where Standard Model-inspired channels
determine whether identity is preserved, transformed, branched, confined, or
annihilated.}
\end{minipage}}
\end{center}

\bigskip

The three state-loss identities:

\begin{align}
  \Delta I_{ij} \leq 0 \quad&\Rightarrow\quad \text{ordinary interaction loss} \\
  L_{ij} = -\Delta I_{ij} \quad&\Rightarrow\quad \text{measured state loss} \\
  \Delta I_{ij} > 0 \quad&\Rightarrow\quad \text{short-lived positive information event}
\end{align}

% ============================================================================
\section{File and Library Reference}
\label{sec:files}
% ============================================================================

\begin{center}
\renewcommand{\arraystretch}{1.25}
\begin{tabular}{ll}
\toprule
\textbf{File} & \textbf{Contents} \\
\midrule
\texttt{IDENTITY/identity\_matrix.hpp}       & Matrix2x2, Matrix3x3, free helpers \\
\texttt{IDENTITY/particle\_identity.hpp}     & ParticleIdentity, SpinProxy, constructors \\
\texttt{IDENTITY/interaction\_channel.hpp}   & Channel row, gates, kappa scoring \\
\texttt{IDENTITY/interaction\_candidate.hpp} & InteractionCandidate, evaluate\_pair() \\
\texttt{IDENTITY/identity.hpp}               & Umbrella include \\
\texttt{PARTICLES/particles.hpp}             & Extended: includes IDENTITY/identity.hpp \\
\texttt{COLLISION/qcd\_force.hpp}            & Extended: qcd\_accumulate\_forces\_sm() \\
\bottomrule
\end{tabular}
\end{center}

\vspace{8pt}
\noindent\textbf{Dependency chain (header-only, single direction):}

\begin{center}
\texttt{identity\_matrix}
$\;\longrightarrow\;$ \texttt{particle\_identity}
$\;\longrightarrow\;$ \texttt{interaction\_channel}
$\;\longrightarrow\;$ \texttt{interaction\_candidate}
\end{center}

\vspace{8pt}
\noindent\textbf{IKK theory documents (companion reading):}
\begin{itemize}
  \item \texttt{IKK\_Doc\_I\_Schrodinger\_Alternative} ---
    scale-ladder definition, identity bundle, hidden-active condition,
    projection failure, entanglement relation-particle.
  \item \texttt{IKK\_Doc\_II\_Initial\_Draft} ---
    scale axioms, matrix packing chain, entropy and complexity,
    dark matter kernel equations, Python and C++ prototypes,
    open variable register.
  \item \texttt{IKK\_Doc\_III\_Kernel\_V4\_Projection} ---
    bilinear projection form for radial density.
  \item \texttt{section\_ikk\_kernel\_bridge} ---
    formal mapping of IKK theory to VSIM \texttt{IDENTITY/} kernel:
    \texttt{ScaleLevel}, \texttt{scale\_mask}, \texttt{fuzz\_ratio()},
    \texttt{H\_charge/color/spin}, hidden-active flags, open variable register.
\end{itemize}

\end{document}
